\RequirePackage[T1]{fontenc}
\documentclass[10pt,conference]{IEEEtran}
\usepackage{amsmath,amssymb,amsthm,graphicx,booktabs,array,url,flushend}
\usepackage[hidelinks]{hyperref}
\usepackage{tikz}
\usetikzlibrary{arrows.meta,positioning}
\emergencystretch=2em
\newtheorem{theorem}{Theorem}
\newtheorem{lemma}{Lemma}
\newtheorem{proposition}{Proposition}
\theoremstyle{definition}
\newtheorem{definition}{Definition}
\newtheorem{remarknum}{Remark}
\newcommand{\code}[1]{\texttt{#1}}
\newcommand{\G}[1]{G#1}
\newcommand{\tool}{ReleaseWeave}
\newcommand{\lang}{RIL}
\begin{document}
\title{Recovering Versioned Releases Across Build and Publish Boundaries}
\maketitle
\begin{abstract}
Releasing interdependent packages across build systems and registries can
leave consumers unfinished after their providers succeed. We present a
saga-inspired release protocol combining deterministic planning, stage-aware
scheduling, and consumer-owned release obligations. Git history and release
records form the ledger without an additional persistent release-state store.
Planning is a pure function of fixed history, records, graph, and configuration.
Consumer windows retain undelivered intent after provider success, while
separate train and fresh-work windows distinguish prerelease targets from
completed work. Under an explicit retry invariant, the single-history model
establishes eight conditional guarantees covering recovery, version stability,
ledger discharge, convergence, and scope/channel behavior. All 120 controlled
fault episodes recover after repair; an analytical source-discharge alternative
loses unfinished work in 84 (70\%). This compares admission rules, not competing
tools or production failure rates. An accompanying agent guide supports plan
inspection and CI release invocation. Proposed fleet and worker extensions
preserve source-owned records and reuse verified provider outputs. A prospective
LLVM scenario projects a Linux release-time reduction from an upstream
154-minute baseline to about 103 minutes (33\%) on 6--10 compatible 16-core
machines, assuming a wider graph, sharded tests, and five-minute output transfers;
no distributed speedup has been measured. Full implementation refinement and
external exactly-once effects remain outside the claims.
\end{abstract}
\begin{IEEEkeywords}
software release engineering, continuous integration and delivery,
dependency scheduling, recovery, idempotence
\end{IEEEkeywords}
\input{introduction}
\input{system}
\input{model}
\input{guarantees}
\input{evaluation}
\input{discussion}
\clearpage
\IEEEtriggeratref{34}
\input{references}
\end{document}
